\section{Discussion}
\label{sec:discussion}

\paragraph{Modality complementarity.}
The asymmetry between AUC and downstream Sharpe across modalities is a key finding.
The price Transformer achieves competitive predictive AUC but near-zero downstream Sharpe on the synthetic corpus, suggesting that price patterns are highly contested at 15-minute resolution.
Web intelligence, by contrast, contributes an additional directional cue not captured by price patterns alone, consistent with the hypothesis that exogenous event information can matter at intraday timescales.
The new \textsc{PriceWeb} comparison in Table~\ref{tab:main} sharpens this point: structured web scalars alone recover only a partial gain over price-only, while remaining materially below \CGCMA{} in mean Sharpe.
The real-news main comparison (Section~\ref{sec:main}) is consistent with this complementarity under more diverse sentiment conditions, but it should still be read as evidence from a short stress-test horizon rather than as proof of cross-regime stability.

\paragraph{Asynchronous alignment and the absorption window.}
Our bar-aligned lag analysis reveals a non-monotonic pattern: web intelligence utility peaks in the 30--60 min post-publication window (signal Sharpe $+0.57$), weakens but remains positive at 60--90 min ($+0.28$), and becomes strongly negative beyond 90 min ($-1.33$).
Very fresh news ($<$30 min) shows near-zero utility ($-0.30$), which is consistent with the possibility that the immediate market response is already underway before the next prediction bar.
We treat this \emph{absorption-window} pattern as an empirical regularity that motivates staleness-aware fusion, not as a definitive causal account of market microstructure.

\paragraph{Conditional gating vs.\ exponential staleness decay.}
\SACMA{} uses an exponential schedule with a single global decay constant $\lambda$, which applies identically to all assets, all sources, and all quality levels.
\CGCMA{} replaces this with a learned per-dimension gate conditioned on price--text agreement, web signal quality, and modality lag.
Table~\ref{tab:design} sharpens the comparison against simpler lag-aware controls.
Fixed decay, hard stale filtering, and the dual-branch predecessor all improve on the price-only baseline, which is consistent with lag being a useful conditioning signal.
At the same time, none of these simpler controls recovers the full mean gain of \CGCMA{}, and the learned scalar decay control remains near parity.
The mechanism is straightforward: when text is low-quality or stale in a given fold, the gate network can suppress the cross-modal context and fall back toward the price CLS prediction, whereas simpler controls apply the same decay or filtering rule without access to the full grounded context.
The dedicated component sweep in Section~\ref{sec:design} further supports this interpretation: removing lag, cross-attention, or vector-valued gating each lowers mean Sharpe, even though the no-gate residual variant is not strictly monotonic and therefore leaves one residual design question open.
Table~\ref{tab:main} still shows that this learned per-dimension gating approach yields the highest mean downstream Sharpe across the evaluated models.
At the same time, the run-level and matched-fold intervals reported in Section~\ref{sec:results} remain wide, so the most defensible claim at the current sample size is one of favorable ranking and effect direction rather than definitive separation from every baseline.

\paragraph{Deep vs.\ linear models and the role of text diversity.}
A key finding from the synthetic corpus is that deep cross-modal attention models require not just sufficient sample scale, but also sufficient text signal diversity.
At ${\sim}500$ training samples per fold with 99\% bullish synthetic text, \SACMA{}-LG degrades from $+0.158$ to $+0.021$---a reversal of the multimodal advantage despite adequate scale.
The real-news corpus restores positive multimodal Sharpe across several fusion models, providing empirical support for the text diversity hypothesis; only \textsc{TFN} remains negative, consistent with its over-parameterization risk at this evaluation scale.
This finding has a practical implication: multimodal systems should validate that their text training signal is directionally balanced; periods of near-monotone sentiment can hurt fusion performance.
The current paper should therefore be read as a controlled study of asynchronous multimodal fusion against representative generic architectures and a small set of task-specific lag-aware controls under a shared protocol, not as an exhaustive benchmark over all possible alternatives.
